\documentclass[conference]{IEEEtran}
\IEEEoverridecommandlockouts

% ================== PAQUETES ==================
\usepackage[utf8]{inputenc}
\usepackage[spanish,es-tabla]{babel}
\usepackage{amsmath,amssymb,amsfonts}
\usepackage{graphicx}
\usepackage{placeins}
\usepackage{textcomp}
\usepackage{xcolor}
\usepackage{cite}
\usepackage{caption}
\usepackage{subcaption}
\usepackage{float}
\usepackage{hyperref}
\usepackage{fancyhdr}
\usepackage{booktabs}
\usepackage{tabularx}

% ===== Ajustes de espaciado IEEE =====
\interdisplaylinepenalty=2500
\setlength{\textfloatsep}{6pt plus 1pt minus 2pt}
\setlength{\floatsep}{6pt plus 1pt minus 2pt}
\setlength{\intextsep}{6pt plus 1pt minus 2pt}
\setlength{\abovedisplayskip}{4pt}
\setlength{\belowdisplayskip}{4pt}

% ================= ENCABEZADO =================
\pagestyle{fancy}
\fancyhf{}
\fancyhead[L]{\includegraphics[height=0.8cm]{logo_umng.png}}
\fancyhead[C]{\small\textit{Control Lineal y Laboratorio - UMNG}}
\fancyhead[R]{\small\thepage}
\renewcommand{\headrulewidth}{0.4pt}

\begin{document}
\raggedbottom

\title{Modelado de Sistemas Mecatrónicos\\
{\normalsize Informe de Laboratorio}}

\author{
\IEEEauthorblockN{Daniel García Araque}
\IEEEauthorblockA{\textit{Ingeniería Mecatrónica}\\
\textit{Universidad Militar Nueva Granada}}
\and
\IEEEauthorblockN{David Santiago García Suarez}
\IEEEauthorblockA{\textit{Ingeniería Mecatrónica}\\
\textit{Universidad Militar Nueva Granada}}
}

\maketitle
\thispagestyle{fancy}

\begin{center}
\textbf{Asignatura:} Control Lineal y Laboratorio\\
\textbf{Programa:} Ingeniería en Mecatrónica
\end{center}

%====================================================
\begin{abstract}
El presente laboratorio aborda el modelado de sistemas mecatrónicos mediante el método de Newton-Euler y Euler-Lagrange. Se desarrollan ecuaciones diferenciales, funciones de transferencia y representaciones en espacio de estado para distintos sistemas físicos, incluyendo modelos lineales y no lineales.
\end{abstract}

\begin{IEEEkeywords}
Modelado mecatrónico, sistemas lineales, espacio de estados, funciones de transferencia.
\end{IEEEkeywords}

%====================================================
\section{Introducción}

El modelado matemático constituye una etapa esencial en el análisis y diseño de sistemas de control. En sistemas mecatrónicos, donde interactúan dominios mecánicos, eléctricos e hidráulicos, resulta necesario obtener representaciones matemáticas que permitan estudiar la dinámica y facilitar el diseño de controladores.

Las metodologías de Newton-Euler y Euler-Lagrange proporcionan herramientas robustas para describir la evolución en el tiempo de sistemas complejos y permiten obtener modelos equivalentes en diferentes dominios matemáticos.

%====================================================
\section{Marco Teórico}

El modelado consiste en describir la relación entre entradas, estados y salidas mediante leyes físicas fundamentales.

\begin{table}[htbp]
\centering
\caption{Analogías Físicas entre Sistemas Mecatrónicos}
\footnotesize
\begin{tabularx}{\columnwidth}{lXXX}
\toprule
\textbf{Dominio} & \textbf{Esfuerzo} & \textbf{Flujo} & \textbf{Elemento} \\
\midrule
Eléctrico & Voltaje & Corriente & R, L, C \\
Mecánico & Fuerza/Torque & Velocidad & Masa, resorte, amortiguador \\
Hidráulico & Presión & Caudal & Resistencia hidráulica \\
Térmico & Temperatura & Flujo térmico & Capacidad térmica \\
\bottomrule
\end{tabularx}
\end{table}

%====================================================
\section{Desarrollo de la Práctica}

%================ SISTEMA 1 ========================
\subsection{Sistema 1: Sistema Electrohidráulico}

\begin{figure}[htbp]
\centering
\includegraphics[width=0.9\columnwidth]{Sistema_ElectroH}
\caption{Sistema electrohidráulico utilizado para el modelado y análisis dinámico.}
\label{fig:sistema_electrohidraulico}
\end{figure}

\subsubsection{Modelado matemático del sistema}

El sistema electrohidráulico está compuesto por un amplificador, un motor de corriente continua, una válvula y un tanque con comportamiento hidráulico. Se describen a continuación las cuatro variables de estado: la corriente de armadura $i_a$, la velocidad angular $\omega$, la posición angular de la válvula $\theta$ y el nivel del fluido $h$.

Aplicando las leyes de Kirchhoff, Newton y el balance de masa, se obtiene el modelo dinámico no lineal:

\begin{equation}
\dot{i}_a = \frac{1}{L_a}(k_a v - R_a i_a - K_b\omega)
\end{equation}

\begin{equation}
\dot{\omega} = \frac{K_m}{J}i_a
\end{equation}

\begin{equation}
\dot{\theta} = \omega
\end{equation}

\begin{equation}
\dot{h} = \frac{K_v}{A}\theta - \frac{k}{A}\sqrt{h}
\end{equation}

donde el término $K_b\omega$ representa la fuerza contraelectromotriz generada por el movimiento del motor.

\subsubsection{Linealización del Sistema}

El sistema se linealiza alrededor del punto de operación $\bar{H} = 1.5$ m. Utilizando una expansión de Taylor de primer orden para el término no lineal:

\begin{equation}
k\sqrt{h} \approx k\sqrt{\bar{H}} + \frac{k}{2\sqrt{\bar{H}}}\Delta h
\end{equation}

La ecuación linealizada del tanque queda expresada como:

\begin{equation}
A\Delta\dot{h} = K_v\Delta\theta - \frac{k}{2\sqrt{\bar{H}}}\Delta h
\end{equation}

\subsubsection{Funciones de Transferencia}

Aplicando la transformada de Laplace con condiciones iniciales nulas, se obtienen las siguientes funciones de transferencia:

\textbf{Sistema electro-mecánico:}
\begin{equation}
G_1(s) = \frac{\Theta(s)}{V(s)} = \frac{k_a K_m}{s(L_a Js^2 + R_a Js + K_b K_m)}
\end{equation}

\textbf{Sistema hidráulico:}
\begin{equation}
G_2(s) = \frac{H(s)}{\Theta(s)} = \frac{K_v/A}{s + \frac{k}{2A\sqrt{\bar{H}}}}
\end{equation}

\textbf{Función de transferencia total:}
\begin{equation}
G(s) = G_1(s) \cdot G_2(s)
\end{equation}

\subsubsection{Cálculo de la Constante de la Válvula}

Para que el sistema hidráulico mantenga el equilibrio con un tiempo de asentamiento $t_s = 5$ s, se calcula:

\begin{equation}
\frac{k}{2A\sqrt{\bar{H}}} = \frac{1}{5}
\end{equation}

Despejando:
\begin{equation}
k = 0.2 \cdot 2 \cdot 50 \cdot \sqrt{1.5} = 24.495
\end{equation}

\subsubsection{Simulación y Análisis de Resultados}

Se diseñó un diagrama en Simulink con dos ramas paralelas: una rama que implementa el modelo no lineal completo del sistema y otra que implementa el modelo linealizado. Ambas ramas reciben la misma entrada de voltaje y permiten comparar directamente sus respuestas.

\textbf{Implementación:}
\begin{itemize}
\item \textbf{Modelo No Lineal:} Se implementaron las cuatro ecuaciones diferenciales del sistema, manteniendo el término $\sqrt{h}$ en la ecuación del tanque.
\item \textbf{Modelo Lineal:} Se utilizó el modelo linealizado alrededor de $\bar{H} = 1.5$ m, reemplazando el término no lineal por su aproximación de primer orden.
\item \textbf{Entrada:} Escalón de voltaje $v(t) = 5$ V aplicado en $t = 0$ s.
\end{itemize}

\textbf{Observaciones a partir de la simulación:}
\begin{itemize}
\item El nivel del tanque $h(t)$ debe aumentar gradualmente desde el estado inicial hasta alcanzar un nuevo punto de equilibrio.
\item Para pequeñas variaciones alrededor de $\bar{H}$, ambos modelos (lineal y no lineal) deben presentar respuestas muy similares, validando la linealización.
\item El modelo no lineal presenta un comportamiento asintótico más realista debido al término $\sqrt{h}$ en el flujo de salida.
\item La velocidad angular $\omega(t)$ alcanza un valor estacionario proporcional al voltaje aplicado.
\item La corriente de armadura $i_a(t)$ presenta un transitorio inicial y luego se estabiliza en un valor constante.
\item El tiempo de asentamiento del nivel hidráulico debe ser aproximadamente 5 segundos, según el diseño de la constante $k$.
\item Para grandes variaciones del nivel, se observan diferencias significativas entre el modelo lineal y no lineal, evidenciando las limitaciones de la linealización.
\end{itemize}

\begin{figure}[h]
    \centering
    \includegraphics[width=0.5\textwidth]{scope1.png}
    \caption{Scope Función de transferencia}
    \label{fig:mi_imagen}
\end{figure}


\begin{figure}[h]
    \centering
    \includegraphics[width=0.5\textwidth]{matlab1.png}
    \caption{Diagrama con Matlab Function}
    \label{fig:mi_imagen}
\end{figure}




\FloatBarrier

%================ SISTEMA 2 ========================
\subsection{Sistema 2: Articulación Flexible}

\begin{figure}[htbp]
\centering
\includegraphics[width=0.9\columnwidth]{Sistema_Arti}
\caption{Sistema de articulación flexible considerado para el modelado rotacional.}
\label{fig:sistema_articulacion}
\end{figure}

A continuación se presenta el desarrollo del modelo matemático para el sistema de articulación flexible mostrado en la Figura~\ref{fig:sistema_articulacion}.

\subsubsection{Modelo Matemático en Ecuaciones de Primer Orden}

Escribiendo las ecuaciones de Newton-Euler para rotación para las dos articulaciones obtenemos las ecuaciones:

\begin{equation}
J_1\ddot{\phi}_1(t) = T_m(t) - k(\phi_1(t) - \phi_2(t)) - b(\dot{\phi}_1(t) - \dot{\phi}_2(t))
\end{equation}

\begin{equation}
J_2\ddot{\phi}_2(t) = k(\phi_1(t) - \phi_2(t)) + b(\dot{\phi}_1(t) - \dot{\phi}_2(t))
\end{equation}

Donde el torque del motor es $T_m(t) = K_t i(t)$.

\textbf{Ecuación Eléctrica:} Aplicando la ley de voltajes de Kirchhoff en la malla del motor y considerando el circuito como uno electromecánico:

\begin{equation}
V(t) = Ri(t) + L\frac{di(t)}{dt} + e_b(t)
\end{equation}

Donde la fuerza contraelectromotriz es $e_b(t) = K_e\dot{\phi}_1(t)$.

\subsubsection{Representación en Espacio de Estados (Primer Orden)}

Se definen las variables de estado: $x_1 = \phi_1, \; x_2 = \dot{\phi}_1, \; x_3 = \phi_2, \; x_4 = \dot{\phi}_2, \; x_5 = i$.

De las ecuaciones dinámicas y cinemáticas realizadas, se obtiene:

\begin{equation}
\dot{x}_1 = x_2
\end{equation}

\begin{equation}
\dot{x}_2 = \frac{1}{J_1}[K_t x_5 - k(x_1 - x_3) - b(x_2 - x_4)]
\end{equation}

\begin{equation}
\dot{x}_3 = x_4
\end{equation}

\begin{equation}
\dot{x}_4 = \frac{1}{J_2}[k(x_1 - x_3) + b(x_2 - x_4)]
\end{equation}

\begin{equation}
\dot{x}_5 = \frac{1}{L}[V(t) - Rx_5 - K_e x_2]
\end{equation}

\subsubsection{Funciones de Transferencia}

Aplicando la transformada de Laplace a las ecuaciones dinámicas con condiciones iniciales nulas:

\textbf{1. Función $G_2(s) = \frac{\Phi_2(s)}{\Phi_1(s)}$:} Relaciona el movimiento de la carga con el eje del motor:

\begin{equation}
G_2(s) = \frac{\Phi_2(s)}{\Phi_1(s)} = \frac{bs + k}{J_2s^2 + bs + k}
\end{equation}

\textbf{2. Función $G_3(s) = \frac{\Phi_1(s)}{I(s)}$:} Relaciona la posición del codo con la corriente de entrada (considerando la carga acoplada):

\begin{equation}
G_3(s) = \frac{\Phi_1(s)}{I(s)} = \frac{K_t(J_2s^2 + bs + k)}{s[J_1J_2s^2 + b(J_1 + J_2)s + k(J_1 + J_2)]}
\end{equation}

\textbf{3. Función $G_1(s) = \frac{I(s)}{V(s)}$:} Incorporado la dinámica eléctrica:

\begin{equation}
G_1(s) = \frac{I(s)}{V(s)} = \frac{G_2(s)}{Ls + R + K_e G_3(s)}
\end{equation}

\textbf{4. Función $G_4(s) = \frac{\Phi_1(s)}{V(s)}$:} Por regla de la cadena entre las funciones anteriores:

\begin{equation}
G_4(s) = G_1(s) \cdot G_2(s) \cdot G_3(s)
\end{equation}

\subsubsection{Asignación de Parámetros}

Se busca que $G_2(s)$ tenga respuesta \emph{subamortiguada} ($\zeta < 1$) y $G_4(s)$ tenga respuesta \emph{sobreamortiguada} ($\zeta > 1$). Analizando los denominadores característicos:

\begin{itemize}
\item Para $G_2(s)$: $\zeta_2 = \frac{b}{2\sqrt{kJ_2}}$. Se requiere $\zeta_2 < 1$.
\item Para $G_4(s)$: $\zeta_4 = \frac{b}{\sqrt{kJ_{eq}}}$, donde $J_{eq} = \frac{J_1J_2}{J_1+J_2}$. Se requiere $\zeta_4 > 1$.
\end{itemize}

Debido a la relación de inercia, se seleccionan los siguientes parámetros para cumplir las condiciones:

\begin{table}[htbp]
\centering
\caption{Parámetros del Sistema de Articulación Flexible}
\begin{tabular}{ccc}
\toprule
\textbf{Parámetro} & \textbf{Valor} & \textbf{Unidad} \\
\midrule
$J_1$ & $2.0$ & kg·m$^2$ \\
$J_2$ & $1.0$ & kg·m$^2$ \\
$k$ & $50$ & Nm \\
$b$ & $25$ & Nm·s \\
$K_t$ & $0.3$ & - \\
$K_e$ & $0.3$ & - \\
\bottomrule
\end{tabular}
\end{table}

\subsubsection{Simulación y Análisis de Resultados}

Se diseñó un diagrama en Simulink con dos ramas paralelas: una rama que utiliza el modelo matricial completo en espacio de estados (primer orden) y otra rama que utiliza bloques de funciones de transferencia para validar la equivalencia entre ambas representaciones.

\textbf{Implementación:}
\begin{itemize}
\item \textbf{Modelo en Espacio de Estados:} Se implementó el sistema de 5 ecuaciones diferenciales de primer orden usando matrices $A$, $B$, $C$ y $D$.
\item \textbf{Modelo con Funciones de Transferencia:} Se implementó $G_4(s)$ utilizando bloques en cascada y retroalimentación.
\item \textbf{Entrada:} Escalón de voltaje $V(t) = 12$ V aplicado en $t = 0$ s.
\end{itemize}

\begin{figure}[h]
    \centering
    \includegraphics[width=0.5\textwidth]{scope2.png}
    \caption{Diagrama con Matlab Function}
    \label{fig:mi_imagen}
\end{figure}

\textbf{Observaciones esperadas:}
\begin{itemize}
\item La posición angular del eje motor $\phi_1(t)$ debe presentar una respuesta sobreamortiguada sin oscilaciones, alcanzando su valor final de manera suave y gradual.
\item La posición angular del eje de carga $\phi_2(t)$ debe mostrar una respuesta subamortiguada con oscilaciones amortiguadas características de un sistema de segundo orden con $\zeta < 1$.
\item Se debe observar un desfase temporal entre $\phi_1(t)$ y $\phi_2(t)$, evidenciando la flexibilidad mecánica del acoplamiento.
\item La corriente de armadura $i(t)$ presenta un pico inicial seguido de una disminución hasta alcanzar un valor estacionario relacionado con la carga mecánica.
\item Las oscilaciones en $\phi_2(t)$ se deben a la interacción entre la rigidez $k$ y la inercia $J_2$, con amortiguamiento proporcionado por $b$.
\item Ambas implementaciones (espacio de estados y funciones de transferencia) deben generar respuestas idénticas, validando la equivalencia matemática de las representaciones.
\item El sistema exhibe dos modos dinámicos dominantes: uno lento asociado al movimiento global y uno rápido asociado a las vibraciones mecánicas.
\item La diferencia entre $\phi_1$ y $\phi_2$ representa la deformación del elemento flexible, la cual oscila inicialmente antes de estabilizarse.
\end{itemize}

\FloatBarrier

\begin{figure}[h]
    \centering
    \includegraphics[width=0.5\textwidth]{doscope.png}
    \caption{Scope 3 funcion de transferencia}
    \label{fig:mi_imagen}
\end{figure}

%================ SISTEMA 3 ========================
\subsection{Sistema 3: Sistema Mecánico mediante Euler-Lagrange}

\begin{figure}[htbp]
\centering
\includegraphics[width=0.9\columnwidth]{Sistema_Mec}
\caption{Sistema mecánico empleado para el modelado mediante Euler-Lagrange.}
\label{fig:sistema_mecanico}
\end{figure}

Para obtener las ecuaciones de movimiento, se utiliza el método energético de Euler-Lagrange.

\subsubsection{Coordenadas Generalizadas}

Se define el vector de coordenadas generalizadas $q(t) \in \mathbb{R}^3$ como:

\begin{equation}
q = \begin{bmatrix} x(t) \\ r(t) \\ \theta(t) \end{bmatrix}
\end{equation}

Donde:
\begin{itemize}
\item $x$: Posición horizontal de la masa $m_1$ (carro).
\item $r$: Longitud instantánea del resorte.
\item $\theta$: Ángulo del péndulo con la vertical.
\end{itemize}

\subsubsection{Cinemática}

Se establece un marco de referencia inercial en la posición inicial de $m_1$. Las posiciones son:

\textbf{Para la masa $m_1$:}
\begin{equation}
x_1 = x
\end{equation}
\begin{equation}
y_1 = 0
\end{equation}

\textbf{Para la masa $m_2$ (Péndulo):}
\begin{equation}
x_2 = x + r\sin\theta
\end{equation}
\begin{equation}
y_2 = -r\cos\theta
\end{equation}

Derivando respecto al tiempo obtenemos las velocidades. La velocidad al cuadrado de $m_2$ es necesaria para la energía cinética:

\begin{equation}
\dot{x}_2^2 = \dot{x}^2 + \dot{r}^2 + r^2\dot{\theta}^2 + 2\dot{x}\dot{r}\sin\theta + 2\dot{x}r\dot{\theta}\cos\theta
\end{equation}

\subsubsection{Energías y Lagrangiano}

La energía cinética ($T$) total del sistema es:

\begin{equation}
\begin{split}
T = \frac{1}{2}(m_1 + m_2)\dot{x}^2 + \frac{1}{2}m_2(\dot{r}^2 + r^2\dot{\theta}^2) \\
+ m_2\dot{x}\dot{r}\sin\theta + m_2\dot{x}r\dot{\theta}\cos\theta
\end{split}
\end{equation}

La energía potencial ($V$) del sistema, considerando el resorte y la gravedad (con referencia en $y = 0$), es:

\begin{equation}
V = \frac{1}{2}k(r - l_0)^2 - m_2gr\cos\theta
\end{equation}

El Lagrangiano $L = T - V$ resulta en:

\begin{equation}
\begin{split}
L = \frac{1}{2}(m_1 + m_2)\dot{x}^2 + \frac{1}{2}m_2(\dot{r}^2 + r^2\dot{\theta}^2) \\
+ m_2\dot{x}r\dot{\theta}\cos\theta + m_2\dot{x}\dot{r}\sin\theta \\
- \frac{1}{2}k(r - l_0)^2 + m_2gr\cos\theta
\end{split}
\end{equation}

\subsubsection{Ecuaciones de Movimiento}

Aplicando $\frac{d}{dt}\left(\frac{\partial L}{\partial \dot{q}_i}\right) - \frac{\partial L}{\partial q_i} = 0$:

\textbf{Dinámica de la coordenada $x$:}
\begin{equation}
(m_1 + m_2)\ddot{x} + m_2\ddot{r}\sin\theta + m_2r\ddot{\theta}\cos\theta = m_2r\dot{\theta}^2\sin\theta - 2m_2\dot{r}\dot{\theta}\cos\theta
\end{equation}

\textbf{Dinámica de la coordenada $r$:}
\begin{equation}
m_2\sin\theta\,\ddot{x} + m_2\ddot{r} = m_2r\dot{\theta}^2 - k(r - l_0) + m_2g\cos\theta
\end{equation}

\textbf{Dinámica de la coordenada $\theta$:}
\begin{equation}
m_2r\cos\theta\,\ddot{x} + m_2r^2\ddot{\theta} = -2m_2r\dot{r}\dot{\theta} - m_2gr\sin\theta
\end{equation}

\subsubsection{Linealización del Modelo}

El sistema se linealiza alrededor del punto de equilibrio estático $(x = 0, \; r = r_{eq}, \; \theta = 0)$.

\subsubsection{Punto de Equilibrio ($r_{eq}$)}

De la ecuación dinámica para $r$ en reposo:
\begin{equation}
0 = 0 - k(r_{eq} - l_0) + m_2g(1)
\end{equation}

Despejando $r_{eq}$:
\begin{equation}
r_{eq} = l_0 + \frac{m_2g}{k}
\end{equation}

\subsubsection{Modelo de Pequeña Señal}

Asumiendo $\sin\theta \approx \theta$, $\cos\theta \approx 1$, $\dot{\theta}^2 \approx 0$, $\theta \approx 0$, las ecuaciones se simplifican como:

\textbf{Resorte (Desacoplado):} La ecuación de $r$ se convierte en:
\begin{equation}
m_2\ddot{r} + k\delta r = 0
\end{equation}

Esto indica que, en el modelo lineal, el resorte se comporta como un oscilador armónico simple independientemente del péndulo.

\textbf{Péndulo y Carro (Acoplados):} Las ecuaciones para $x$ y $\theta$ resultan en:

\begin{equation}
(m_1 + m_2)\ddot{x} + m_2r_{eq}\ddot{\theta} = 0
\end{equation}

\begin{equation}
\ddot{x} + r_{eq}\ddot{\theta} = -g\theta
\end{equation}

Resolviendo este sistema algebraico para despejar las aceleraciones:

\begin{equation}
\ddot{x} = \frac{m_2g}{m_1}\theta
\end{equation}

\begin{equation}
\ddot{\theta} = -\frac{(m_1 + m_2)g}{m_1r_{eq}}\theta
\end{equation}

\subsubsection{Parámetros del Sistema}

Se utilizaron los siguientes valores para la simulación en Matlab/Simulink:

\begin{table}[htbp]
\centering
\caption{Parámetros del Sistema Mecánico}
\begin{tabular}{ccc}
\toprule
\textbf{Parámetro} & \textbf{Valor} & \textbf{Unidad} \\
\midrule
$m_1$ & $2.0$ & kg \\
$m_2$ & $1.0$ & kg \\
$k$ & $50$ & N/m \\
$l_0$ & $0.5$ & m \\
$g$ & $9.81$ & m/s$^2$ \\
\bottomrule
\end{tabular}
\end{table}
\begin{figure}[h]
    \centering
    \includegraphics[width=0.5\textwidth]{s.png}
    \caption{Diagrama con Matlab Function}
    \label{fig:mi_imagen}
\end{figure}
\subsubsection{Simulación y Análisis de Resultados}

Se diseñó un diagrama en Simulink con dos ramas: una implementando el modelo no lineal completo mediante las ecuaciones de Euler-Lagrange y otra implementando el modelo linealizado de pequeña señal.

\textbf{Implementación:}
\begin{itemize}
\item \textbf{Modelo No Lineal:} Se programaron las tres ecuaciones diferenciales acopladas de segundo orden, manteniendo los términos $\sin\theta$, $\cos\theta$ y $\dot{\theta}^2$.
\item \textbf{Modelo Lineal:} Se utilizaron las ecuaciones linealizadas con aproximaciones $\sin\theta \approx \theta$ y $\cos\theta \approx 1$.
\item \textbf{Condiciones Iniciales:} Se aplicaron pequeñas perturbaciones en $\theta_0$ y $r_0$ para observar la dinámica del sistema.
\item \textbf{Fuerza Externa:} Se consideró $F_x = 0$ (sistema libre sin fuerza aplicada en el carro).
\end{itemize}

\textbf{Observaciones esperadas:}
\begin{itemize}
\item \textbf{Desacoplamiento del resorte:} En el modelo lineal, la coordenada $r$ oscila de manera independiente como un oscilador armónico simple con frecuencia natural $\omega_r = \sqrt{k/m_2}$, sin interactuar con las otras coordenadas.
\item \textbf{Acoplamiento carro-péndulo:} Las coordenadas $x$ y $\theta$ están fuertemente acopladas. Cuando se perturba $\theta$, el carro se desplaza horizontalmente y viceversa.
\item \textbf{Inestabilidad del péndulo invertido:} Si el péndulo parte de una posición cercana a $\theta = 0$ (arriba), el modelo lineal muestra crecimiento exponencial de $\theta(t)$, característico de un péndulo invertido inestable.
\item \textbf{Diferencias entre modelos:} Para pequeñas oscilaciones iniciales ($|\theta_0| < 10°$), ambos modelos presentan respuestas similares. Para ángulos mayores, el modelo no lineal muestra desviaciones significativas.
\item \textbf{Conservación de energía:} En ausencia de amortiguamiento, el modelo no lineal debe conservar la energía total del sistema, oscilando indefinidamente.
\item \textbf{Efectos no lineales:} El modelo completo captura fenómenos como la dependencia de la frecuencia de oscilación con la amplitud inicial (oscilaciones anarmónicas).
\item \textbf{Comportamiento del carro:} El desplazamiento $x(t)$ del carro presenta oscilaciones debido al acoplamiento con el péndulo. En el modelo lineal, esta oscilación crece sin límite si el sistema es inestable.
\item \textbf{Validación de la linealización:} La linealización es válida únicamente en un vecindario pequeño del punto de equilibrio; fuera de esta región, los modelos divergen considerablemente.
\end{itemize}

\FloatBarrier

%================ SISTEMA 4 ========================
\subsection{Sistema 4: Controladores Analógicos PD, PI y PID}

\begin{figure}[htbp]
\centering
\includegraphics[width=0.9\columnwidth]{Sistema_Circuito}
\caption{Implementación analógica de controladores PD y PI.}
\label{fig:sistema_circuito}
\end{figure}

\subsubsection{Modelado de controladores}

El análisis se realiza mediante impedancias equivalentes en amplificadores operacionales ideales.

\textbf{Controlador PD (Proporcional-Derivativo):}
\begin{equation}
G_{PD}(s)=\frac{R_4R_2}{R_3R_1}(1+R_1C_1s)
\end{equation}

donde la ganancia proporcional es $K_p = \frac{R_4R_2}{R_3R_1}$ y el tiempo derivativo es $T_d = R_1C_1$.

\textbf{Controlador PI (Proporcional-Integral):}
\begin{equation}
G_{PI}(s)=\frac{R_4R_2}{R_3R_1}
\left(1+\frac{1}{R_2C_2s}\right)
\end{equation}

donde la ganancia proporcional es $K_p = \frac{R_4R_2}{R_3R_1}$ y el tiempo integral es $T_i = R_2C_2$.

\textbf{Controlador PID (Proporcional-Integral-Derivativo):}

La acción PID combina las tres acciones de control y se expresa en el dominio del tiempo como:

\begin{equation}
u(t)=K_pe(t)+\frac{K_p}{T_i}\int e(t)dt+K_pT_d\frac{de(t)}{dt}
\end{equation}

En el dominio de Laplace:
\begin{equation}
G_{PID}(s) = K_p\left(1 + \frac{1}{T_is} + T_ds\right)
\end{equation}

\subsubsection{Parámetros de Diseño}

Para la implementación analógica, se seleccionaron los siguientes valores de componentes:

\begin{table}[htbp]
\centering
\caption{Valores de Componentes para Controladores Analógicos}
\begin{tabular}{ccc}
\toprule
\textbf{Componente} & \textbf{Valor} & \textbf{Controlador} \\
\midrule
$R_1$ & $10$ k$\Omega$ & PD, PI, PID \\
$R_2$ & $10$ k$\Omega$ & PD, PI, PID \\
$R_3$ & $10$ k$\Omega$ & PD, PI, PID \\
$R_4$ & $10$ k$\Omega$ & PD, PI, PID \\
$C_1$ & $10$ $\mu$F & PD, PID \\
$C_2$ & $10$ $\mu$F & PI, PID \\
\bottomrule
\end{tabular}
\end{table}

\subsubsection{Simulación y Análisis de Resultados}

Se implementaron los tres controladores analógicos (PD, PI y PID) en Simulink utilizando bloques de funciones de transferencia que representan los circuitos con amplificadores operacionales. Cada controlador se conectó en lazo cerrado con una planta de prueba.

\textbf{Planta de Prueba:}

Se utilizó una función de transferencia de segundo orden como planta genérica:
\begin{equation}
G_p(s) = \frac{\omega_n^2}{s^2 + 2\zeta\omega_ns + \omega_n^2}
\end{equation}

con $\omega_n = 5$ rad/s y $\zeta = 0.3$ (subamortiguado).

\textbf{Implementación:}
\begin{itemize}
\item \textbf{Sistema sin Controlador:} Respuesta en lazo abierto de la planta.
\item \textbf{Control PD:} Implementación de la acción proporcional-derivativa.
\item \textbf{Control PI:} Implementación de la acción proporcional-integral.
\item \textbf{Control PID:} Implementación combinada de las tres acciones.
\item \textbf{Entrada:} Referencia tipo escalón unitario $r(t) = 1$ aplicado en $t = 0$ s.
\end{itemize}

\begin{figure}[h]
    \centering
    \includegraphics[width=0.5\textwidth]{4.png}
    \caption{Diagrama con Matlab Function}
    \label{fig:mi_imagen}
\end{figure}

\begin{figure}[h]
    \centering
    \includegraphics[width=0.5\textwidth]{4s.png}
    \caption{Scope 4}
    \label{fig:mi_imagen}
\end{figure}


\textbf{Observaciones:}

\textbf{Sistema sin control:}
\begin{itemize}
\item Respuesta subamortiguada con oscilaciones significativas.
\item Error en estado estacionario nulo para entrada escalón (sistema tipo 0).
\item Tiempo de establecimiento prolongado debido al bajo amortiguamiento.
\end{itemize}

\textbf{Control PD:}
\begin{itemize}
\item Reducción notable del sobreimpulso ($M_p$) comparado con el sistema sin control.
\item Disminución del tiempo de establecimiento ($t_s$) debido al incremento del amortiguamiento efectivo.
\item La acción derivativa anticipa cambios en el error, mejorando la estabilidad relativa.
\item \textbf{Limitación importante:} No elimina el error en estado estacionario para perturbaciones o cambios en la referencia tipo rampa.
\item Sensibilidad al ruido de alta frecuencia en la señal de error debido al término derivativo.
\end{itemize}

\textbf{Control PI:}
\begin{itemize}
\item Eliminación completa del error en estado estacionario gracias a la acción integral.
\item La acción integral acumula el error a lo largo del tiempo, generando una señal de control que compensa desviaciones permanentes.
\item Posible incremento del sobreimpulso y del tiempo de establecimiento comparado con el sistema sin control.
\item Mejora significativa en el rechazo a perturbaciones de tipo escalón.
\item \textbf{Riesgo de windup:} Si la señal de control se satura, el integrador puede acumular error excesivo.
\end{itemize}

\textbf{Control PID:}
\begin{itemize}
\item Combina las ventajas de PD y PI: elimina error estacionario (integral) y mejora velocidad de respuesta (derivativo).
\item Respuesta más rápida con menor sobreimpulso que el control PI puro.
\item Mayor robustez ante perturbaciones comparada con PD o PI individualmente.
\item Permite alcanzar un compromiso óptimo entre rapidez, precisión y estabilidad mediante el ajuste adecuado de $K_p$, $T_i$ y $T_d$.
\item El ajuste de parámetros es más complejo, requiriendo técnicas como Ziegler-Nichols o métodos analíticos.
\item Sensibilidad moderada al ruido debido al término derivativo, aunque menor que en PD puro si se filtra adecuadamente.
\end{itemize}

\textbf{Comparación entre controladores:}
\begin{itemize}
\item PD es superior para mejorar la velocidad de respuesta y reducir oscilaciones, pero no elimina errores estacionarios.
\item PI es esencial cuando se requiere precisión en estado estacionario, aunque puede sacrificar velocidad.
\item PID ofrece el mejor desempeño global cuando los parámetros están bien sintonizados.
\item La elección depende de las especificaciones del sistema: si se requiere error nulo en estado estacionario, se debe incluir acción integral (PI o PID).
\end{itemize}

\FloatBarrier


\section{Análisis de Resultados}

Las simulaciones permitieron validar la coherencia entre los modelos matemáticos y la dinámica observada. La linealización es válida alrededor del punto de operación, mientras que los modelos no lineales describen fenómenos adicionales.

%====================================================
\section{Conclusiones}

\begin{enumerate}
\item Se obtuvieron modelos matemáticos coherentes para diferentes sistemas mecatrónicos.
\item La linealización simplifica el análisis pero reduce fidelidad fuera del equilibrio.
\item Newton-Euler y Euler-Lagrange son metodologías complementarias.
\end{enumerate}

%====================================================
\begin{thebibliography}{00}
\bibitem{ogata} K. Ogata, \emph{Ingeniería de Control Moderna}, Pearson, 2010.
\bibitem{khalil} H. Khalil, \emph{Nonlinear Systems}, Prentice Hall, 2002.
\end{thebibliography}

\end{document}